\section{Methods}
\label{app:methods}

This appendix records the algebra behind the results: the bias-overbound
reconstruction, the positive-semidefiniteness of the bias cross-covariance, the
correlated-$\ge$-independent inequality, the Cholesky whitening and the
common-mode statistic, and the finite-ceiling blind-spot bound with its
mutual-exclusivity argument. All are consistent with the committed modules of
Section~\ref{sec:availability}.

\subsection{The bias-overbound reconstruction}
\label{app:overbound}

A metrology laboratory reports a Type-B expanded uncertainty $U_i$ at coverage
factor $k_{\mathrm{cov}}$, so the standard uncertainty is $u_i=U_i/k_{\mathrm{cov}}$
\citep{jcgm100}. To carry the systematic component as a one-sided bounded bias at
an allocated integrity tail rather than as a Gaussian noise sigma
(Remark~\ref{rem:category}), we scale $u_i$ by the two-sided normal quantile at that
tail and add the deterministic ageing allowance,
\begin{equation}
  b_i \;=\; u_i\,\Phi^{-1}\!\Big(1-\tfrac{\mathrm{tail}}{2}\Big)+\mathrm{ageing}_i
       \;=\; \frac{U_i}{k_{\mathrm{cov}}}\,\Phi^{-1}\!\Big(1-\tfrac{\mathrm{tail}}{2}\Big)
             +\mathrm{ageing}_i,
  \tag{\ref{eq:overbound}}
\end{equation}
which is~\eqref{eq:overbound}. The multiplier $\Phi^{-1}(1-\mathrm{tail}/2)$ realizes
a zero-mean Gaussian CDF-overbound (DeCleene-style paired overbounding) of the
traceability-bias error: it bounds $|X|\le b_i$ at two-sided probability
$1-\mathrm{tail}$, so the tail scaling is exact under that overbound, which is the
operative modelling assumption we own here rather than a claim of physical
Gaussianity (Remark~\ref{rem:category}). The deep-tail choice---the demonstration
uses a two-sided $\mathrm{tail}=\num{2e-7}$, i.e.\ a per-side $\num{e-7}$ giving
$\Phi^{-1}(1-\num{e-7})$---is \Modelled{}, because no monthly bulletin has sample
support at that probability (Section~\ref{sec:tracescope}). The ageing term is added
outside the quantile scaling because it is deterministic, so its worst case over the
exposure interval is its endpoint value and it is not inflated by the tail factor.

\subsection{Positive semidefiniteness of $\Sib{}$}
\label{app:psd}

With $D=\operatorname{diag}(b_1,\dots,b_N)$ and correlation matrix
$R=(\rho_{ij})$ (symmetric, unit diagonal, $R\succeq0$), the bias
cross-covariance~\eqref{eq:sigmab} is the congruence
\begin{equation}
  \Sib{} \;=\; D\,R\,D,
  \label{eq:congruence}
\end{equation}
since $(DRD)_{ij}=b_i\rho_{ij}b_j$. For any $v\in\R^N$, set $w=Dv$; then
$v^{\!\top}\Sib{}v=w^{\!\top}Rw\ge0$ because $R\succeq0$. Hence
$\Sib{}\succeq0$ (Proposition~\ref{prop:psd}). A block-constant $R$ built from
disjoint shared-realizer groups---each group an all-$\rho$ block, zero across
groups---is positive semidefinite for $\rho\in[0,1]$ (each block has eigenvalues
$1+(m-1)\rho\ge0$ and $1-\rho\ge0$ for a group of size $m$), so the demonstration's
pattern is admissible.

\subsection{Correlated fusion dominates independent fusion}
\label{app:corr}

Expanding the quadratic form with $w_i\ge0$, $b_i\ge0$, $\rho_{ij}\in[0,1]$,
\begin{equation}
  w^{\!\top}\Sib{}\,w
  \;=\; \sum_i w_i^2 b_i^2 \;+\; \sum_{i\neq j} \rho_{ij}\,w_i w_j\,b_i b_j
  \;\ge\; \sum_i w_i^2 b_i^2,
  \label{eq:expand}
\end{equation}
because every off-diagonal term is nonnegative; taking square roots (both sides
nonnegative, $\sqrt{\cdot}$ monotone) gives~\eqref{eq:corrgeqindep}. The gap is
$\sum_{i\neq j}\rho_{ij}w_iw_jb_ib_j$, nonzero only for shared-realizer pairs
($\rho_{ij}>0$), and strictly positive whenever two such sources carry nonzero
weight (Theorem~\ref{thm:corr}). In the fully coherent limit $\rho_{ij}\equiv1$
within a group, $\sum_{i,j\in\text{grp}}w_iw_jb_ib_j=(\sum_{i\in\text{grp}}w_ib_i)^2$,
so the correlated form is the squared within-group \emph{linear} sum. This is the
key to caveat~I3: the companion's linear fusion $\sum_i w_i b_i$ satisfies
\begin{equation}
  \Big(\sum_i w_i b_i\Big)^2
  = \sum_i w_i^2 b_i^2 + \sum_{i\neq j} w_i w_j b_i b_j
  \;\ge\; \sum_i w_i^2 b_i^2 + \sum_{i\neq j}\rho_{ij} w_i w_j b_i b_j
  = w^{\!\top}\Sib{}w,
  \label{eq:lindominates}
\end{equation}
since $\rho_{ij}\le1$; the companion protection level's bias term therefore
dominates the correlated fusion. Equality holds \emph{iff} every weight-carrying
pair lies in a single shared realizer at $\rho_{ij}=1$ (all weighted sources in one
realizer group); with two or more distinct realizer groups the linear
across-all-sources form \emph{strictly} dominates $w^{\!\top}\Sib{}w$ \emph{even at}
$\rho\equiv1$, since $\Sib{}$ then omits the cross-group products $w_iw_jb_ib_j$ that
the linear square retains. The correlated-bias theorem thus warns against a
variance-style allocation and does not harden the linear companion PL
(Remark~\ref{rem:i3}); the strict multi-group dominance only reinforces that it
cannot.

\subsection{Cholesky whitening and the common-mode statistic}
\label{app:cmstat}

With residual covariance $\Om{}\succ0$, the Cholesky factor $L$ ($\Om{}=LL^{\!\top}$,
$L$ lower-triangular, positive diagonal) whitens $z=L^{-1}r$ with
$\operatorname{Cov}(z)=L^{-1}\Om{}L^{-\top}=I_N$ \citep{aitken1935}. The
generalized-least-squares estimate of a common shift $c$ (model $r=\bmone c+\text{noise}$,
$\operatorname{Cov}=\Om{}$) minimizes $(r-\bmone c)^{\!\top}\Om^{-1}(r-\bmone c)$,
giving $\hat c=(\bmone^{\!\top}\Om^{-1}\bmone)^{-1}\bmone^{\!\top}\Om^{-1}r$ with
$\operatorname{Var}(\hat c)=(\bmone^{\!\top}\Om^{-1}\bmone)^{-1}$. The
variance-normalized square is
\begin{equation}
  \scm{} \;=\; \frac{\hat c^{\,2}}{\operatorname{Var}(\hat c)}
         \;=\; \frac{(\bmone^{\!\top}\Om^{-1}r)^2}{\bmone^{\!\top}\Om^{-1}\bmone},
  \tag{\ref{eq:cmstat}}
\end{equation}
which is~\eqref{eq:cmstat}; under Gaussian residuals and the null $c=0$ it is the
square of a unit-variance normal, hence $\scm{}\sim\chi^2_1$. For the post-fit
residual $r=y-\bmone\hat x$ with
$\hat x=(\bmone^{\!\top}\Om^{-1}\bmone)^{-1}\bmone^{\!\top}\Om^{-1}y$, one has
$\bmone^{\!\top}\Om^{-1}r=\bmone^{\!\top}\Om^{-1}y-(\bmone^{\!\top}\Om^{-1}\bmone)\hat
x=0$, so $\scm{}\equiv0$ identically~\eqref{eq:postfitzero}: the fit consumes the
common-mode degree of freedom, and the statistic is informative only on a pre-fit
residual against an independent anchor (caveat~I2, Remark~\ref{rem:i2}).

\subsection{The finite-ceiling blind-spot bound}
\label{app:blind}

Let a unit fault direction $d$ enter the pre-fit residual, $r=d\,\xi$. Whiten:
$L^{-1}d\,\xi$. The idealized solution-separation (parity) detector projects onto
the orthogonal complement of the whitened geometry $L^{-1}\bmone{}$; with
$P=I-\tfrac{(L^{-1}\bmone)(L^{-1}\bmone)^{\!\top}}{\|L^{-1}\bmone\|^2}$ its statistic
is $\propto\|P L^{-1}d\|\,\xi$, so its trip magnitude is
$\alpha_{\mathrm{ss}}(d)\propto\|P L^{-1}d\|^{-1}$, infinite iff $P L^{-1}d=0$, i.e.\
$L^{-1}d\in\operatorname{span}(L^{-1}\bmone)$, i.e.\ $d\parallel\bmone{}$. The
common-mode detector~\eqref{eq:cmstat} has statistic
$\propto|\bmone^{\!\top}\Om^{-1}d|\,\xi/\sqrt{\bmone^{\!\top}\Om^{-1}\bmone}$, so its
trip magnitude is $\alpha_{\mathrm{cm}}(d)\propto|\bmone^{\!\top}\Om^{-1}d|^{-1}$,
infinite iff $\bmone^{\!\top}\Om^{-1}d=0$. These are the two blind
conditions~\eqref{eq:blindconds}, and the ceiling is
$\xi^\star(d)=\min(\alpha_{\mathrm{ss}},\alpha_{\mathrm{cm}})$~\eqref{eq:ceiling}.

\paragraph{Mutual exclusivity under a PD whitening (caveat~I1).}
Suppose $\Om{}\succ0$, so $\Om^{-1}\succ0$ and in particular
$\bmone^{\!\top}\Om^{-1}\bmone{}>0$. If the first blind condition holds,
$d=\lambda\bmone{}$ with $\lambda\neq0$, then
$\bmone^{\!\top}\Om^{-1}d=\lambda\,\bmone^{\!\top}\Om^{-1}\bmone{}\neq0$, so the
second fails and $\alpha_{\mathrm{cm}}$ is finite. If the second holds,
$\bmone^{\!\top}\Om^{-1}d=0$, then $d$ cannot equal any nonzero $\lambda\bmone{}$
(which would give the strictly positive form above), so $d\not\parallel\bmone{}$,
$P L^{-1}d\neq0$, and $\alpha_{\mathrm{ss}}$ is finite. The two blind sets meet only
at $d=0$; for every nonzero $d$ at least one detector has a finite trip magnitude,
so $\xi^\star(d)<\infty$ (Theorem~\ref{thm:finite}). The ceiling is infinite only
for $d=0$ or a rank-deficient (non-PD) $\Om{}$. Finally, the idealized parity
$\alpha_{\mathrm{ss}}$ omits the false-alert threshold and Bonferroni allocation of a
real MHSS subset detector, which raise the real trip magnitude; the modelled
$\alpha_{\mathrm{ss}}$ is therefore smaller than the operational one, and
$\xi^\star(d)$ is a conservative \emph{lower} bound on the true undetectable
magnitude.
